\section{État de l'art}
\label{etat-de-lart}
\shorthandoff{;:!?}

Ce chapitre poursuit un objectif unique, subordonné à la finalité du mémoire~:
identifier, parmi les approches existantes, celles qui permettent d'\textbf{établir une
mesure fiable de la charge de travail des équipes Loan Servicing et d'en déduire le
coût d'un deal --- rétrospectivement, pour un deal traité, et de façon prédictive, pour
un deal futur}. La charge de travail n'est pas étudiée pour elle-même~: elle est la
grandeur intermédiaire qui, mesurée puis valorisée, donne accès au coût, puis à la
rentabilité. Cet objectif recoupe celui assigné à la mission~: «~estimer avec précision
le coût associé à un nouveau deal~» et obtenir «~un coût par événement et par deal pour
chaque équipe LS~».

Plutôt que de juxtaposer les disciplines, nous les ordonnons le long d'un fil unique ---
de la mesure du temps au coût --- et n'évaluons chaque méthode que par sa contribution à
ce fil, selon trois questions~: mesure-t-elle le temps par événement~? résiste-t-elle à
un back office nearshore, automatisé et hétérogène~? conduit-elle à un coût imputable au
deal~? Le périmètre se resserre ainsi sur deux finalités articulées par un socle
commun --- le temps unitaire par événement~: \textbf{mesurer} le coût d'un deal traité
(§\ref{sec:abc}--\ref{sec:mesure-temps}), et le \textbf{prédire} pour un deal futur
(§\ref{sec:prediction}).

% =============================================================================
\subsection{Du coût des activités au coût par événement~: ABC et TDABC}
\label{sec:abc}
% =============================================================================

\subsubsection{Le principe de l'Activity-Based Costing}

Toute méthode de coût doit répondre à une question~: comment répartir des charges
\emph{indirectes} --- salaires d'un back office, systèmes, encadrement --- sur les
objets qui les consomment~? La comptabilité traditionnelle répond par des clés
grossières~: on répartit au prorata d'un volume ou d'heures de main-d'œuvre. Le défaut
est connu~: un deal simple à fort volume et un deal complexe à faible volume reçoivent
une part de coût sans rapport avec les ressources qu'ils exigent réellement.
L'\textit{Activity-Based Costing} (ABC), formalisé par Kaplan et Cooper à la fin des
années 1980, corrige ce biais en interposant les \textbf{activités}~: on trace d'abord
le coût des ressources vers les activités qu'elles servent (\textit{resource drivers}),
puis le coût des activités vers les objets qui les déclenchent (\textit{activity
drivers}) \cite{kaplan1988}. La littérature distingue à ce titre deux natures
d'inducteurs~: les \textit{transaction drivers}, qui comptent le nombre de fois qu'une
activité survient, et les \textit{duration drivers}, qui mesurent le temps qu'elle
consomme --- ces derniers étant plus précis lorsque les occurrences d'une même activité
n'ont pas toutes la même durée. En banque, l'ABC a montré qu'il révèle des clients ou
produits dont le service consomme un back office disproportionné par rapport aux revenus
qu'ils génèrent --- exactement notre problème de rentabilité par deal.

\paragraph{Pertinence et limite pour notre contexte.}
Le principe de l'ABC est directement pertinent~: notre «~objet~» est le deal, nos
«~activités~» sont les événements Loan Servicing, et nous cherchons un
\textit{duration driver} --- le temps par événement. Mais l'ABC classique a une faiblesse
qui le disqualifie en l'état~: il estime le temps par activité au moyen d'\textbf{enquêtes
déclaratives}, où chaque opérateur indique le pourcentage de son temps consacré à chaque
tâche. Ces enquêtes sont coûteuses à mener, subjectives, et surtout se périment dès que
les processus changent --- ce qui, sous déploiement continu de RPA, est notre quotidien.
De plus, l'ABC suppose implicitement que la capacité est utilisée à 100~\%, et noie donc
le coût de la capacité inutilisée dans des taux gonflés. Ces deux défauts --- coût de
maintenance et cécité à la capacité oisive --- sont précisément ceux que la variante
suivante corrige.

\subsubsection{Le Time-Driven ABC~: deux paramètres, une équation}

Kaplan et Anderson introduisent en 2004 le \textit{Time-Driven Activity-Based Costing}
(TDABC), qui réduit toute la mécanique de l'ABC à \textbf{deux paramètres} par groupe de
ressources \cite{kaplan2004}. Le premier est le \textbf{taux de coût de capacité}
(\textit{capacity cost rate})~: le coût total du groupe (salaires, encadrement,
systèmes, locaux) divisé par sa \emph{capacité pratique} --- non pas le temps
théoriquement disponible, mais celui réellement travaillé, estimé à environ 80--85~\% du
théorique une fois retranchés pauses, formations et aléas. On obtient un coût par minute.
Le second est un ensemble d'\textbf{équations de temps} (\textit{time equations}) qui
estiment la durée d'une transaction à partir de ses caractéristiques~: un temps de base,
augmenté d'incréments pour chaque facteur qui l'allonge. La forme canonique est~:

\begin{equation}
t = \beta_0 + \sum_{k} \beta_k \, x_k
\label{eq:time-equation}
\end{equation}

\noindent où $\beta_0$ est le temps de base, $x_k$ les caractéristiques de la transaction
(par exemple~: contrôle manuel requis, expédition à l'export) et $\beta_k$ le temps
additionnel qu'elles imposent. Le coût de l'objet est alors la somme, sur tous ses
événements, du temps consommé multiplié par le taux de coût de capacité~:

\begin{equation}
C = \sum_{i=1}^{N} t_i \times c_u
\label{eq:tdabc}
\end{equation}

Cette équation est, à la notation près, \textbf{celle que notre analyse terrain a
établie} --- $C = (T_{\text{LIQ}} + T_{\text{AC}}) \times N \times w_s$ (chapitre~1). Le
TDABC est donc le socle de notre finalité rétrospective. Il présente deux vertus
supplémentaires décisives. D'une part, il rend \emph{visible la capacité inutilisée}~:
comme on a valorisé la capacité \emph{fournie}, la différence entre minutes disponibles
et minutes réellement consommées apparaît comme une ligne propre --- information cruciale
pour dimensionner les équipes FR/PT. D'autre part, l'équation~\ref{eq:tdabc} est
\emph{nativement prédictive}~: connaître les temps unitaires permet de simuler le coût
d'un volume ou d'un mix futurs, ce qui fait le pont vers notre seconde finalité.

\paragraph{La limite critique --- et c'est notre gap.}
La littérature postérieure, moins promotionnelle que les textes fondateurs, pointe une
faiblesse unique mais décisive~: \emph{la précision du TDABC dépend entièrement de la
qualité des estimations de temps}, et la méthode s'applique mal au \textit{knowledge
work} non routinier, où le temps d'une tâche est un mauvais proxy de la ressource
consommée. Le cas le plus proche du nôtre est documenté en santé~: Koster \textit{et al.}
(2023) montrent que l'hétérogénéité des patients rend les temps par activité non
ponctuels mais variables, au point de devoir remplacer les estimations uniques par une
logique floue capturant un intervalle du meilleur au pire cas \cite{koster2023}. La
transposition est immédiate~: nos deals sont à nos équipes ce que les patients
multimorbides sont aux soignants --- une population hétérogène pour laquelle un temps
unitaire figé ne suffit pas. C'est exactement ce que la mission nomme «~\textit{wide
possibilities of post operation process which make the workload difficult to
predict}~». Estimer ce temps unitaire fiable, en environnement hétérogène et automatisé,
est le verrou que ce mémoire doit lever --- et qui justifie d'examiner les méthodes de
mesure du temps.

% =============================================================================
\subsection{Mesurer le temps par événement}
\label{sec:mesure-temps}
% =============================================================================

\subsubsection{L'étude des temps (Taylor)~: chronométrer et normaliser}

La méthode fondatrice de mesure du travail est l'\textit{étude des temps}, formalisée par
Taylor (1911) \cite{taylor1911, groover2007}. Son principe~: décomposer une tâche en
éléments observables, chronométrer chacun sur plusieurs cycles pour obtenir un
\emph{temps observé} $T_o$, puis le corriger en deux étapes. On applique d'abord un
\textbf{facteur d'allure} $F_a$ --- jugement porté par l'analyste sur le rythme de
l'opérateur par rapport à une cadence «~normale~» (un opérateur rapide se voit affecter
$F_a > 1$) --- pour obtenir un \emph{temps normal} $T_n = T_o \times F_a$. On ajoute
ensuite des \textbf{majorations} $A$ pour les aléas incompressibles (fatigue, pauses,
interruptions), d'où le \emph{temps standard}~:

\begin{equation}
T_s = T_o \times F_a \times (1 + A)
\label{eq:temps-standard}
\end{equation}

Ben-Gal \textit{et al.} (2010) ont établi que ce principe reste valide dans les
services, à condition d'adapter la granularité de décomposition aux tâches informatisées
\cite{bengal2010}.

\paragraph{Pertinence et limites pour notre contexte.}
Le principe est le bon~: il produit exactement le $t_i$ qu'exige le TDABC, et l'idée de
décomposer un traitement en événements chronométrables fonde notre \textit{clocking}.
Mais le \emph{protocole} tayloriste échoue sur trois points précis dans un back office
nearshore. \emph{Premièrement}, le facteur d'allure $F_a$ repose sur l'observation
directe d'un analyste~: dans des équipes réparties entre Paris et Lisbonne, cette
observation est rarement praticable, et l'auto-déclaration réintroduit un biais de
désirabilité sociale. \emph{Deuxièmement}, l'étude des temps suppose une «~méthode
standard~» stabilisée~; or l'automatisation redéfinit la façon d'exécuter un même
événement d'un trimestre à l'autre, rendant tout temps standard rapidement obsolète.
\emph{Troisièmement}, elle suppose des tâches nettement délimitées, alors que le travail
réel est entrelacé --- plusieurs deals menés de front, interrompus par des requêtes
\textit{ad hoc} --- ce qui brouille la mesure du temps effectif. Conclusion~: on garde le
principe (chronométrer par événement), on abandonne le protocole (observation +
jugement d'allure) au profit d'une mesure ancrée dans les traces système.

\begin{figure}[htbp]\centering
\begin{tikzpicture}[>=Stealth,node distance=6mm,
  box/.style={rectangle,rounded corners,draw=black!60,fill=blue!8,text width=26mm,minimum height=16mm,align=center,font=\small}]
  \node[box](n1){1. Décomposition de la tâche en événements};
  \node[box,right=of n1](n2){2. Chronométrage sur plusieurs cycles ($T_o$)};
  \node[box,right=of n2](n3){3. Facteur d'allure ($F_a\to T_n$)};
  \node[box,right=of n3](n4){4. Majorations aléas ($T_s$)};
  \draw[->](n1)--(n2);\draw[->](n2)--(n3);\draw[->](n3)--(n4);
\end{tikzpicture}
\caption{Protocole de l'étude des temps (Taylor, 1911). Les étapes 2 et 3 — observation directe et jugement d'allure — sont celles qui résistent le moins à un back office distribué.}
\label{fig:taylor}\end{figure}



\subsubsection{Le MTM~: des temps prédéterminés, mais pour des gestes}

Frank et Lillian Gilbreth déplacent l'attention de la durée vers le \emph{mouvement}~:
ils décomposent tout geste en unités élémentaires (les \textit{therbligs} --- saisir,
positionner, relâcher). De là naît le \textit{Methods-Time Measurement} (MTM), formalisé
par Maynard \textit{et al.} (1948) \cite{gilbreth1917, maynard1948}~: un système de temps
\emph{prédéterminés} qui attribue \textit{a priori} une valeur temporelle normalisée à
chaque mouvement élémentaire, si bien qu'on peut estimer la durée d'une opération
\emph{avant} de l'exécuter, par simple addition, sans chronométrer.

\paragraph{Pertinence et limites pour notre contexte.}
L'attrait est réel~: le MTM promettrait d'estimer $t_i$ sans campagne de mesure. Mais il
présuppose que le travail se réduit à des gestes physiques catalogables --- vrai à
l'atelier, faux au back office. Les «~mouvements~» d'un opérateur Loan Servicing sont des
séquences \emph{cognitives et numériques}~: naviguer entre écrans Loan~IQ, interpréter
une clause, arbitrer entre instructions contradictoires. Leur durée ne dépend pas d'une
mécanique gestuelle mais du \emph{contenu informationnel} du deal~: un tirage bilatéral
en euros et un tirage syndiqué multidevise à vingt participants sollicitent les mêmes
écrans pour des durées radicalement différentes. Aucune table MTM générique ne capture
cette variabilité. Le MTM confirme donc, par son échec, que notre référentiel de temps
doit être \emph{reconstruit empiriquement} à partir de données réelles, non dérivé d'un
catalogue.

\subsubsection{Le Process Mining~: lire le temps dans les traces système}

Le \textit{Process Mining}, formalisé par van der Aalst (2016), résout le problème par un
tout autre angle~: au lieu d'observer l'opérateur, il exploite les \textbf{logs
d'événements} déjà enregistrés par les systèmes d'information \cite{vanderaalst2016}.
Chaque ligne de log fournit au minimum un triplet $(\text{case\_id}, \text{activité},
\text{timestamp})$~: l'identifiant du cas (ici, le deal), l'activité réalisée, et
l'horodatage. En recomposant ces triplets, le Process Mining \emph{découvre} le
processus réellement suivi (et non le processus théorique), mesure les durées entre
activités, et met au jour ce que les cartographies déclaratives masquent~: boucles de
reprise, approbations redondantes, temps d'attente entre étapes. Jans \textit{et al.}
(2013) en démontrent la valeur dans une organisation financière pour l'analyse
multifacette des processus et la détection d'anomalies \cite{jans2013}.

\paragraph{Pertinence et limites pour notre contexte.}
Le Process Mining lève les deux premières failles du chronométrage tayloriste~: il
ne requiert ni observateur ni jugement d'allure (il lit des horodatages objectifs), et
il se met à jour automatiquement à mesure que les processus évoluent --- répondant au
problème de l'obsolescence sous automatisation. Deux limites, toutefois, écartent son
adoption comme méthode de mesure dans notre contexte.

La première est méthodologique, et recoupe précisément la nôtre~: l'horodatage mesure
un \textbf{temps écoulé} entre deux transactions, non un \textbf{temps effectif} de
travail. Entre le \textit{timestamp} d'ouverture et celui de clôture d'un événement dans
Loan~IQ, l'opérateur a pu être interrompu, traiter un autre deal ou attendre une
validation. Le Process Mining ne fournit donc qu'une \emph{borne supérieure bruitée} du
temps unitaire --- celui-là même que nous cherchons à isoler.

La seconde est juridique et éthique. Exploiter les logs revient à traiter des données
qui tracent indirectement l'activité individuelle des opérateurs~: horodatages, cadence,
durée par transaction. Au regard du RGPD et du droit du travail, un tel usage suppose une
finalité explicite et proportionnée, l'information préalable des salariés, la consultation
du CSE et l'interdiction de tout détournement à des fins d'évaluation individuelle ---
autant de conditions qui en alourdissent considérablement l'implémentation.

Pour ces deux raisons --- une mesure qui reste un temps écoulé, et un cadre d'usage lourd
---, nous retenons le \emph{principe} du Process Mining (lire le temps dans les traces
système) sans l'adopter comme méthode. Notre protocole lui préfère un \textit{clocking}
ciblé, qui capture directement le temps effectif sur un échantillon d'événements et
s'affranchit de ces contraintes.


% =============================================================================
\subsection{Prédire le coût d'un deal futur~: méthodes de modélisation}
\label{sec:prediction}
% =============================================================================

Mesurer le coût d'un deal traité résout la première finalité. La seconde --- estimer le
coût d'un deal \emph{avant} son traitement, à partir de ses seules caractéristiques
connues à la signature --- prolonge naturellement le TDABC, dont l'équation de temps
(\ref{eq:time-equation}) est déjà une fonction que l'on peut \emph{apprendre} sur
l'historique. Là où le TDABC calcule un coût \textit{a posteriori} en observant le
volume et l'\textit{aftercare} réellement consommés, un modèle prédictif estime~:

\begin{equation}
\hat{C}_{\text{deal}} = f\!\left(x_1, \dots, x_p\right)
\label{eq:cout-predit}
\end{equation}

\noindent où $f$ est appris sur les deals passés et où les $x_j$ sont les
caractéristiques disponibles dès la signature~: nombre de prêteurs, devise, nombre de
facilités, rôle (agent, participant, bilatéral), produit. Cette section examine les
familles de modèles candidates pour $f$, de la plus interprétable à la plus performante,
en pesant pour chacune le compromis \emph{précision / interprétabilité} --- compromis
central, car un modèle destiné à arbitrer la rentabilité d'un deal devant le Front Office
doit être non seulement juste, mais explicable.

\subsubsection{La régression linéaire multiple~: le modèle de référence}

Le modèle le plus simple postule que le coût (ou le temps) est une combinaison linéaire
des caractéristiques~:

\begin{equation}
Y = \beta_0 + \beta_1 x_1 + \beta_2 x_2 + \cdots + \beta_p x_p + \varepsilon
\label{eq:regression}
\end{equation}

\noindent où $Y$ est la variable à prédire, les $\beta_j$ des coefficients estimés par
moindres carrés, et $\varepsilon$ un terme d'erreur. Sa force est l'\textbf{interprétabilité
par conception}~: chaque coefficient $\beta_j$ se lit directement comme la contribution
marginale d'un facteur --- «~un prêteur supplémentaire ajoute $\beta$ minutes de
traitement~». C'est un argument de poids pour la communication managériale, et cela
répond directement au sous-objectif de la mission~: «~\textit{understand the impact of
each parameter of the LS workload}~».

\paragraph{Hypothèses et pourquoi elles se brisent ici.}
Cette lisibilité repose sur des hypothèses que notre contexte viole. La régression
suppose la \textbf{linéarité} (l'effet d'un facteur est constant, quel que soit son
niveau) et surtout l'\textbf{additivité} (l'effet total est la somme des effets
individuels). Or nos données terrain montrent des \textbf{interactions}~: un deal
\emph{à la fois} syndiqué, multidevise \emph{et} à clauses spéciales génère une charge
supérieure à la somme des trois effets pris isolément --- la complexité se
\emph{multiplie}, elle ne s'additionne pas. La régression sous-estimerait donc
systématiquement les deals les plus complexes, précisément ceux dont le coût importe le
plus. On peut ajouter manuellement des termes d'interaction ($x_1 \times x_2$), mais leur
nombre explose et leur choix devient arbitraire. La régression reste donc un
\emph{étalon} de référence et un outil de communication, mais insuffisante comme modèle
final.

\subsubsection{Les arbres de décision (CART)~: capturer les interactions}

L'arbre de décision de type CART (\textit{Classification and Regression Trees}) adopte
une logique inverse~: au lieu de poser une équation globale, il \textbf{partitionne
récursivement} l'espace des caractéristiques. À chaque nœud, l'algorithme cherche la
question binaire (par exemple «~nombre de prêteurs $> 50$~?~») qui sépare le mieux les
données, puis recommence sur chaque sous-groupe, jusqu'à des feuilles où la prédiction
est la moyenne des cas qu'y tombent. Pour une cible continue comme le temps, le
«~meilleur~» découpage est celui qui \textbf{minimise la variance} intra-groupe --- on
cherche le seuil qui rend les deux sous-groupes les plus homogènes possible. La
réduction de variance obtenue par un découpage s'écrit~:

\begin{equation}
\Delta = \mathrm{Var}(S) - \frac{|S_L|}{|S|}\,\mathrm{Var}(S_L) - \frac{|S_R|}{|S|}\,\mathrm{Var}(S_R)
\label{eq:variance-reduction}
\end{equation}

\noindent où $S$ est le groupe parent, $S_L$ et $S_R$ les deux sous-groupes issus du
découpage, et $|\cdot|$ leur effectif. L'algorithme retient, à chaque nœud, le découpage
qui maximise $\Delta$.

\paragraph{Pertinence et limites.}
L'arbre a deux vertus majeures pour nous. Il capture \textbf{nativement les
interactions}~: en enchaînant «~si rôle = agent \emph{puis} si multidevise = oui~», il
modélise exactement l'effet multiplicatif que la régression manque, sans qu'on ait à le
spécifier. Et il est \textbf{lisible}~: un chemin de la racine à une feuille se lit comme
une règle métier explicite. Mais un arbre unique souffre de deux défauts. Il est
\textbf{instable} (une variation minime des données change toute sa structure ---
\emph{high variance}), et il ne sait pas \textbf{extrapoler}~: entraîné sur des deals de
2 à 1000 prêteurs, il plafonnera sa prédiction pour un deal à 1500 prêteurs à la valeur
de sa feuille la plus haute. Seul, il sur-apprend et généralise mal~: il faut l'agréger.

\subsubsection{Random Forest et gradient boosting~: agréger pour fiabiliser}

Deux stratégies d'agrégation d'arbres dominent la pratique sur données tabulaires comme
les nôtres. Le \textbf{Random Forest} entraîne des centaines d'arbres \emph{en parallèle},
chacun sur un échantillon aléatoire des données et des variables, puis \emph{moyenne}
leurs prédictions~:

\begin{equation}
\hat{f}(x) = \frac{1}{B}\sum_{b=1}^{B} T_b(x)
\label{eq:random-forest}
\end{equation}

\noindent où $T_b$ est le $b$-ième arbre. En moyennant des arbres décorrélés, la forêt
annule leur instabilité individuelle~: elle est robuste, demande peu de réglage, et
fournit une mesure d'\emph{importance des variables} (via la réduction de variance
cumulée qu'apporte chaque variable) --- utile pour hiérarchiser les inducteurs de charge.
Le \textbf{gradient boosting} (dont XGBoost est l'implémentation de référence) procède au
contraire \emph{séquentiellement}~: chaque nouvel arbre corrige les erreurs résiduelles
des précédents. Cette correction ciblée le rend généralement plus précis que la forêt sur
données tabulaires structurées, au prix d'un réglage plus délicat (un boosting mal réglé
sur-apprend). Le compromis pratique est clair~: Random Forest pour la robustesse et la
simplicité de mise en œuvre~; gradient boosting pour la performance maximale quand on
peut investir dans le réglage.

\paragraph{Le prix à payer~: l'interprétabilité, et comment la restaurer.}
Ces ensembles gomment la lisibilité de l'arbre unique~: on ne peut plus lire une règle,
seulement constater une performance. C'est un vrai problème dans un contexte bancaire, où
une estimation de coût opposée au Front Office doit être justifiable. La littérature sur
l'\textit{interpretable machine learning} répond par des méthodes d'explication
\emph{post-hoc}, appliquées après l'entraînement, indépendamment du modèle. La plus
établie est celle des \textbf{valeurs de Shapley} (SHAP), issue de la théorie des jeux
coopératifs \cite{lundberg2017}. L'idée~: traiter chaque caractéristique comme un
«~joueur~» contribuant à la prédiction, et lui allouer équitablement sa part de l'écart
entre la prédiction et la moyenne. La valeur SHAP $\phi_j$ d'une variable $j$ est sa
contribution marginale moyenne, calculée sur tous les ordres d'entrée possibles des
variables~:

\begin{equation}
\phi_j = \sum_{S \subseteq N \setminus \{j\}} \frac{|S|!\,(|N|-|S|-1)!}{|N|!}\left[f(S \cup \{j\}) - f(S)\right]
\label{eq:shap}
\end{equation}

\noindent où $N$ est l'ensemble des variables et $S$ un sous-ensemble n'incluant pas $j$.
Une propriété fondamentale rend SHAP idéal pour notre usage~: les valeurs SHAP sont
\textbf{additives} et somment exactement à l'écart entre la prédiction du deal et le coût
moyen. On peut donc présenter au Front Office une décomposition intelligible~: «~ce deal
coûte 180~K€, soit 90~K€ au-dessus de la moyenne~; +50~K€ viennent du nombre de prêteurs,
+30~K€ du rôle d'agent, +10~K€ du multidevise~». On obtient ainsi le meilleur des deux
mondes~: la \emph{précision} d'un modèle non linéaire et une \emph{explication} par deal
aussi lisible qu'une régression.

\paragraph{Le pilotage temporel, en appui.}
Enfin, la part de coût liée au \emph{volume} et à sa saisonnalité (pics de
\textit{quarter-end}) relève de l'analyse de séries temporelles. Les modèles classiques
(ARIMA, SARIMAX) y ajoutent des variables exogènes (pipeline commercial, calendrier
réglementaire), mais la littérature récente souligne leurs limites --- linéarité imposée,
horizons fixes --- et leur préfère désormais des ensembles d'arbres et des approches sous
contrainte pour des demandes hétérogènes à horizon variable \cite{park2026}. Ces
approches restent en appui de nos deux finalités~: elles prédisent le \emph{volume} $N$,
non le temps unitaire, qui demeure stable.

\paragraph{Une mise en garde de niveau.}
Deux erreurs guettent l'usage de ces modèles, et structurent notre méthodologie.
D'abord, le modèle prédit le coût \emph{total} d'un deal (fonction du volume et des
caractéristiques), \emph{non} le temps de \textit{booking} unitaire, qui reste stable
autour de vingt minutes~: confondre les deux niveaux serait une faute d'analyse. Ensuite,
et surtout, la qualité de toute prédiction est \textbf{bornée par celle de la mesure}~:
un modèle appris sur des temps unitaires mal estimés prédira un coût faux, quelle que
soit sa sophistication. La seconde finalité est donc entièrement tributaire de la
première --- ce qui justifie l'ordre de ce chapitre et le poids que nous mettons sur la
mesure du temps.

% =============================================================================
\subsection{Les solutions professionnelles~: ce que le marché résout, ce qu'il laisse ouvert}
\label{sec:solutions-marche}
% =============================================================================

La question n'a pas échappé aux éditeurs, et deux familles de solutions éclairent par
contraste notre contribution. Les plateformes de \textbf{Back-Office Workforce
Management} (Verint, et les suites décrites par les intégrateurs du secteur) capturent le
temps de production réel des opérateurs à travers leurs applications pour équilibrer les
charges, suivre les SLA et dimensionner les effectifs \cite{verint}. Elles adressent le
\emph{pilotage} de la charge, avec des gains documentés~; mais elles raisonnent en
\emph{capacité et débit}, non en \emph{coût par objet}~: elles disent si une équipe est
surchargée, jamais ce que coûte un deal donné. À l'inverse, les \textbf{suites de
costing} (CostPerform et le TDABC outillé) calculent un coût par activité et par objet,
mais supposent l'intrant temps déjà résolu --- elles industrialisent l'équation, pas la
mesure du $t_i$.

Un courant professionnel récent, sur les projets d'IA en services financiers, formule la
critique la plus utile à notre positionnement~: la majorité des projets d'automatisation
échouent faute d'avoir mesuré, \emph{au niveau activité}, le travail réel avant
d'automatiser --- la donnée applicative («~six heures dans le système~») étant
inexploitable pour décider, seule la donnée d'activité («~tant de minutes sur telle
exception~») l'étant \cite{summittrails}. Ce constat conforte notre séquençage~: la
mesure fine du temps précède et conditionne toute optimisation. Aucune des deux familles
d'outils ne construit le pont qui nous occupe --- du temps unitaire mesuré jusqu'au coût,
mesuré puis prédit, d'un deal de crédit syndiqué hétérogène.

% =============================================================================
\subsection{Synthèse~: un gap resserré sur l'intrant temps du TDABC}
\label{sec:gap}
% =============================================================================

L'examen de la littérature ne débouche pas sur le constat commode qu'«~aucun auteur n'a
tout combiné~». Il désigne un manque unique et localisé. Le \textbf{TDABC fournit
l'équation juste} (équation~\ref{eq:tdabc}) pour reconstruire \emph{et} prédire un coût
par deal, et notre terrain la confirme. Mais il \textbf{suppose résolu son intrant le plus
difficile}, le temps unitaire $t_i$, que la littérature critique désigne unanimement
comme son talon d'Achille --- d'autant plus en environnement hétérogène, où Koster
\textit{et al.} (2023) ont dû recourir à la logique floue pour absorber la variabilité des
cas \cite{koster2023}. Les choix retenus dans cette synthèse constituent donc la
traduction opérationnelle de l'analyse théorique conduite dans ce chapitre~; leur
validité empirique sera ensuite testée puis confirmée sur le terrain aux chapitres~3 et~4.
Or cet intrant se heurte, en opérations de crédit, à trois obstacles cumulatifs établis au
fil du chapitre~: \emph{(i)} les méthodes classiques (Taylor, MTM) reposent sur une
observation ou des tables inadaptées à un back office nearshore et automatisé~;
\emph{(ii)} le Process Mining lève cet obstacle mais ne livre qu'un temps écoulé bruité,
non effectif~; \emph{(iii)} la constante d'\textit{Event Processing} ($\bar{t} \approx
20$~min) s'applique strictement à l'exécution des actions standard de booking et de saisie
dans Loan~IQ, tandis que toute la complexité cognitive, l'interprétation contractuelle et
les échanges inter-équipes sont isolés dans la composante \textit{Aftercare}~; \emph{(iv)}
la variabilité du coût ne vient donc pas du temps de \textit{booking}, stable, mais du
volume et de l'\textit{aftercare}, que les caractéristiques du deal font varier de façon
non additive~; le modèle de régression de l'\textit{Aftercare} peut incorporer des termes
d'interaction ($x_i \cdot x_j$) ou des traits polynomiaux pour capter ces multiplicateurs
de complexité non additifs tout en demeurant pleinement interprétable.

Conclusion~: nous conservons le principe tayloriste --- décomposer le traitement en
parties chronométrables --- mais en abandonnons le protocole subjectif. Notre
\textit{clocking} s'en inspire directement~: nous chronométrons le temps \textbf{par
parties} du traitement d'un événement --- la création de la notice SR/SAR sur
\textit{Hobart}, l'exécution de l'événement dans Loan~IQ, puis la réconciliation dans
\textit{Ambre} (cf. figure~\ref{fig:decomposition}) --- \textbf{sans appliquer de facteur
d'allure}. Là où Taylor corrige le temps observé par un jugement d'allure porté sur le
rythme de l'opérateur, nous nous en tenons au temps effectif mesuré~: dans un back office
distribué entre Paris et Lisbonne, l'observation directe n'est ni praticable ni
souhaitable, et tout jugement d'allure réintroduirait la subjectivité que la méthode
cherche à écarter.


Notre recherche relève ainsi du \textit{design science} \cite{hevner2004}~: elle produit
un artefact --- un protocole de mesure du temps unitaire et sa conversion en coût par
deal, rétrospectif puis prédictif --- dont la validité s'évalue par sa capacité à résoudre
un problème organisationnel réel~: rendre enfin comparable, deal par deal, le coût back
office et le gain front office.

\shorthandon{;:!?}
\newpage
